\documentclass[UTF8,11pt,a4paper]{ctexart}
\usepackage[margin=24mm]{geometry}
\usepackage{amsmath,amssymb,amsthm,mathtools,bm}
\usepackage{tikz-cd,enumitem,longtable,booktabs}
\usepackage{xurl}
\usepackage[unicode,hidelinks]{hyperref}
\allowdisplaybreaks
\setlength{\emergencystretch}{3em}
\setlength{\parindent}{0pt}
\setlength{\parskip}{0.45em}
\newtheorem*{theorem}{Theorem}
\newtheorem*{lemma}{Lemma}
\newtheorem*{proposition}{Proposition}
\newtheorem*{corollary}{Corollary}
\theoremstyle{definition}\newtheorem*{definition}{Definition}
\theoremstyle{remark}\newtheorem*{remark}{Remark}
\DeclareMathOperator{\Hom}{Hom}
\DeclareMathOperator{\End}{End}
\DeclareMathOperator{\Tr}{Tr}
\DeclareMathOperator{\rank}{rank}
\DeclareMathOperator{\Ind}{Ind}
\DeclareMathOperator{\Res}{Res}
\DeclareMathOperator{\im}{Im}
\DeclareMathOperator{\id}{id}
\newcommand{\R}{\mathbb R}
\newcommand{\C}{\mathbb C}
\newcommand{\Q}{\mathbb Q}
\newcommand{\Z}{\mathbb Z}
\newcommand{\N}{\mathbb N}
\newcommand{\E}{\mathbb E}
\newcommand{\Prob}{\mathbb P}
\newcommand{\eps}{\varepsilon}
\begin{document}
\section*{100｜有限Galois扩张的Hilbert90与范数核}
\textbf{作者／集体来源：}J. S. Milne

\textbf{作品：}Fields and Galois Theory, version 5.00 (June 2021)

\textbf{原文入口：}\url{https://www.jmilne.org/math/CourseNotes/FT500.pdf}

\textbf{原文位置与计题范围：}第5章，印刷页65（5.14--5.15）、69--71（5.21(a)(b)、5.23、5.25）。计一题：有限Galois背景下的乘法一阶上同调消失，附循环扩张的范数核表述。

\textbf{获取日期：}2026-09-28。\quad\textbf{复用依据：}CC BY-NC-SA 4.0；原文第2页。

\textbf{整理归属：}下列题面与证明取自所列人类来源，按注明范围转录、摘编；LaTeX环境、标题、换行及中文说明由助手整理。编辑调整在题内说明。

\section*{作者命题与人类证明}

\subsection*{Supporting result: Dedekind's independence theorem (5.14--5.15)}
\begin{theorem}[Dedekind, 5.14]
Let $F$ be a field and $G$ a group. Every finite set $\{\chi_1,\ldots,\chi_m\}$ of group homomorphisms $G\to F^\times$ is linearly independent over $F$, i.e.,
\[
\sum a_i\chi_i=0\quad\text{(as a function $G\to F$)}
\quad\Longrightarrow\quad a_1=0,\ldots,a_m=0.
\]
\end{theorem}
\begin{proof}
We use induction on $m$. For $m=1$, the statement is obvious. Assume it for $m-1$, and suppose that, for some set $\{\chi_1,\ldots,\chi_m\}$ of homomorphisms $G\to F^\times$ and $a_i\in F$,
\[
a_1\chi_1(x)+a_2\chi_2(x)+\cdots+a_m\chi_m(x)=0\quad\text{for all }x\in G.
\]
We have to show that the $a_i$ are zero. As $\chi_1$ and $\chi_2$ are distinct, they will take distinct values on some $g\in G$. On replacing $x$ with $gx$ in the equation, we find that
\[
a_1\chi_1(g)\chi_1(x)+a_2\chi_2(g)\chi_2(x)+\cdots+a_m\chi_m(g)\chi_m(x)=0
\]
for all $x\in G$. On multiplying the first equation by $\chi_1(g)$ and subtracting it from the second, we obtain the equation
\[
a'_2\chi_2+\cdots+a'_m\chi_m=0,\qquad a'_i=a_i(\chi_i(g)-\chi_1(g)).
\]
The induction hypothesis shows that $a'_i=0$ for $i=2,3,\ldots$. As $\chi_2(g)-\chi_1(g)\ne0$, this implies that $a_2=0$, and so
\[
a_1\chi_1+a_3\chi_3+\cdots+a_m\chi_m=0.
\]
The induction hypothesis now shows that the remaining $a_i$ are also zero.
\end{proof}
\begin{corollary}[5.15]
Let $F$ and $E$ be fields, and let $\sigma_1,\ldots,\sigma_m$ be distinct homomorphisms $F\to E$. Then $\sigma_1,\ldots,\sigma_m$ are linearly independent over $E$.
\end{corollary}
\begin{proof}
Apply the theorem to $\chi_i=\sigma_i|F^\times$.
\end{proof}

\subsection*{Crossed homomorphisms and first cohomology (p. 69)}
Let $G$ be a group. A $G$-module is an abelian group $M$ together with an action of $G$, i.e., a map $G\times M\to M$ such that
\begin{enumerate}[label=(\alph*)]
\item $\sigma(m+m')=\sigma m+\sigma m'$ for all $\sigma\in G$, $m,m'\in M$;
\item $(\sigma\tau)(m)=\sigma(\tau m)$ for all $\sigma,\tau\in G$, $m\in M$;
\item $1_Gm=m$ for all $m\in M$.
\end{enumerate}
Let $M$ be a $G$-module. A crossed homomorphism is a map $f:G\to M$ such that
\[
f(\sigma\tau)=f(\sigma)+\sigma f(\tau)\quad\text{for all }\sigma,\tau\in G.
\]
Note that the condition implies that $f(1)=f(1\cdot1)=f(1)+f(1)$, and so $f(1)=0$.

\paragraph{Example 5.21(a).}
Let $f:G\to M$ be a crossed homomorphism. For any $\sigma\in G$,
\begin{align*}
f(\sigma^2)&=f(\sigma)+\sigma f(\sigma),\\
f(\sigma^3)&=f(\sigma)+\sigma f(\sigma)+\sigma^2 f(\sigma),\\
f(\sigma^n)&=f(\sigma)+\sigma f(\sigma)+\cdots+\sigma^{n-1}f(\sigma).
\end{align*}
Thus, if $G$ is a cyclic group of order $n$ generated by $\sigma$, then a crossed homomorphism $f:G\to M$ is determined by its value, $x$ say, on $\sigma$, and $x$ satisfies the equation
\begin{equation}\label{eq:norm-additive}
x+\sigma x+\cdots+\sigma^{n-1}x=0.
\end{equation}
Moreover, if $x\in M$ satisfies \eqref{eq:norm-additive}, then the formulas
$f(\sigma^i)=x+\sigma x+\cdots+\sigma^{i-1}x$ define a crossed homomorphism $f:G\to M$. Thus, for a finite cyclic group $G=\langle\sigma\rangle$, there is a one-to-one correspondence between crossed homomorphisms and the $x\in M$ satisfying \eqref{eq:norm-additive}, given by $f\leftrightarrow f(\sigma)$.

\paragraph{Example 5.21(b).}
For every $x\in M$, we obtain a crossed homomorphism by putting
\[
f(\sigma)=\sigma x-x,\quad\text{all }\sigma\in G.
\]
Such a crossed homomorphism is said to be principal. The sum and difference of two crossed homomorphisms is again a crossed homomorphism, and the sum and difference of two principal crossed homomorphisms is again principal. Thus we can define
\[
H^1(G,M)=\frac{\{\text{crossed homomorphisms}\}}{\{\text{principal crossed homomorphisms}\}}
\]
(quotient abelian group).

\subsection*{The counted theorem: Hilbert 90 (5.23), with its norm formulation (5.25)}
\begin{theorem}[5.23]
Let $E$ be a Galois extension of $F$ with group $G$; then $H^1(G,E^\times)=0$, i.e., every crossed homomorphism $G\to E^\times$ is principal.
\end{theorem}
\begin{proof}
Let $f$ be a crossed homomorphism $G\to E^\times$. In multiplicative notation, this means that
\[
f(\sigma\tau)=f(\sigma)\cdot\sigma(f(\tau)),\quad\sigma,\tau\in G,
\]
and we have to find a $\gamma\in E^\times$ such that $f(\sigma)=\sigma\gamma/\gamma$ for all $\sigma\in G$. Because the $f(\tau)$ are nonzero, Corollary 5.15 implies that
\[
\sum_{\tau\in G}f(\tau)\tau:E\longrightarrow E
\]
is not the zero map, i.e., there exists an $\alpha\in E$ such that
\[
\beta\mathrel{:=}\sum_{\tau\in G}f(\tau)\tau\alpha\ne0.
\]
But then, for $\sigma\in G$,
\begin{align*}
\sigma\beta
&=\sum_{\tau\in G}\sigma(f(\tau))\cdot\sigma\tau(\alpha)\\
&=\sum_{\tau\in G}f(\sigma)^{-1}f(\sigma\tau)\cdot\sigma\tau(\alpha)\\
&=f(\sigma)^{-1}\sum_{\tau\in G}f(\sigma\tau)\sigma\tau(\alpha),
\end{align*}
which equals $f(\sigma)^{-1}\beta$ because, as $\tau$ runs over $G$, so also does $\sigma\tau$. Therefore,
\[
f(\sigma)=\frac{\beta}{\sigma(\beta)}=\frac{\sigma(\beta^{-1})}{\beta^{-1}}.
\]
\end{proof}
Let $E$ be a Galois extension of $F$ with Galois group $G$. We define the norm of an element $\alpha\in E$ to be
\[
\operatorname{Nm}\alpha=\prod_{\sigma\in G}\sigma\alpha.
\]
For $\tau\in G$, $\tau(\operatorname{Nm}\alpha)=\prod_{\sigma\in G}\tau\sigma\alpha=\operatorname{Nm}\alpha$; and so $\operatorname{Nm}\alpha\in F$. The map $\alpha\mapsto\operatorname{Nm}\alpha:E^\times\to F^\times$ is a homomorphism.
We are interested in determining the kernel of the norm map. Clearly an element of the form $\beta/\sigma\beta$ has norm $1$, and our next result shows that, for cyclic extensions, all elements with norm $1$ are of this form.
\begin{corollary}[Hilbert's Theorem 90, 5.25]
Let $E$ be a finite cyclic extension of $F$, and let $\sigma$ generate $\operatorname{Gal}(E/F)$. Let $\alpha\in E^\times$; if $\operatorname{Nm}_{E/F}\alpha=1$, then $\alpha=\beta/\sigma\beta$ for some $\beta\in E$.
\end{corollary}
\begin{proof}
Let $m=[E:F]$. The condition on $\alpha$ is that $\alpha\cdot\sigma\alpha\cdots\sigma^{m-1}\alpha=1$, and so (see 5.21(a)) there is a crossed homomorphism $f:\langle\sigma\rangle\to E^\times$ with $f(\sigma)=\alpha$. Theorem 5.23 now shows that $f$ is principal, which means that there is a $\beta$ with $f(\sigma)=\beta/\sigma\beta$.
\end{proof}

\section*{整理说明与可修改标注（助手）}
\textbf{标准先修：}有限Galois理论、群作用、线性独立；所需Dedekind独立性证明已附入。

\textbf{本题仍需完成的工作：}由乘法上闭链的相容关系构造带权共轭和，先确保非零，再算出其变换规律，最后将循环范数条件编码为上闭链。

\textbf{取材说明：}保留原文有限Galois语境，不扩展为无限群任意映射的陈述。5.21只摘取(a)(b)及定义段，未纳入留作练习的上同调长正合列。支撑定理不另计题数。原PDF文字层已取得；网页截图接口本次未返回页面图像。

\textbf{$c$：}域扩张中的范数方程与上同调消失

\textbf{$a$：}Galois群作用；交叉同态；特征标线性独立；带权共轭和；范数

\texttt{cross\_theory\_status = yes}

\textbf{具体联系及位置：}原文5.23把上闭链条件变成非零带权和的变换规律，得到主上闭链；5.25再把范数为一的元素转换为该类上闭链。

这些标注是非穷尽初稿；未列出不表示未使用。
\end{document}
